\chapter{Jensen--Wiener Correlation Kernel and the Direct RH Frontier}
\label{chap:g024-jensen-wiener}

\status{Exact reparametrization, exact external/internal Wronskian identities, and globally proved first-order complete monotonicity. Higher complete-monotonicity orders remain open. The Riemann Hypothesis remains open.}

\section{Purpose}

The preceding PF$_3$ generations isolate coefficient-level obstructions and local sufficient conditions. This chapter returns to an RH-equivalent correlation-kernel criterion and records a direct sufficient route based on complete monotonicity. It also separates two different $y$-tilts which agree at $y=0$ but are not the same family away from the real axis.

The external criterion used here is Dimitrov--Xu, \texttt{arXiv:1606.05011}. The $y=0$ positive-definite-kernel frontier is consistent with the framework studied by Csordas, \texttt{arXiv:1309.0055}.

\section{Full-line Fourier normalization}

The repository kernel satisfies
\[
 \xi\!\left(\frac12+z\right)=\int_0^\infty \Phi(t)\cosh(zt)\,dt.
\]
Define
\[
 \boxed{K(t)=\frac12\Phi(|t|).}
\]
Then
\[
 \boxed{\Xi(x):=\xi\!\left(\frac12+ix\right)=\int_{\mathbb R}K(t)e^{-ixt}\,dt.}
\]

\section{Dimitrov--Xu order-two correlation}

For a suitable even positive kernel define
\[
 \nu_2(t)=\int_{\mathbb R}(t-2s)^2K(t-s)K(s)\,ds
\]
and
\[
 \Psi_y(t)=\cosh(ty)\nu_2(t).
\]
Dimitrov and Xu prove, for the Riemann kernel,
\[
 \boxed{RH\iff\mathcal T(\Psi_y)\text{ is dense in }L^1(\mathbb R)\quad\text{for every }0<|y|<\frac12.}
\]
By Wiener's $L^1$ Tauberian theorem this is equivalent to absence of real zeros of $\widehat\Psi_y$ for every such $y$.

\section{Centered normal form}

Set
\[
 t=2u,\qquad s=u-r.
\]
Then
\[
 t-s=u+r,\qquad s=u-r,\qquad t-2s=2r,
\]
so
\[
 \boxed{\nu_2(2u)=4\int_{\mathbb R}r^2K(u+r)K(u-r)\,dr.}
\]
Define
\[
 \boxed{C(u):=\int_{\mathbb R}r^2K(u+r)K(u-r)\,dr}
\]
and
\[
 \boxed{D_y(u):=\cosh(2yu)C(u).}
\]
Hence
\[
 \boxed{\Psi_y(2u)=4D_y(u).}
\]
Scaling and fixed dilation preserve the density question. Therefore
\[
 \boxed{RH\iff\mathcal T(D_y)\text{ is dense in }L^1(\mathbb R)\quad\forall\,0<|y|<\frac12.}
\]
Since $D_y$ is real, even and integrable and $\widehat D_y(0)>0$, the Wiener condition can be written
\[
 \boxed{\widehat D_y(x)>0\quad\forall x\in\mathbb R,\quad 0<|y|<\frac12.}
\]
This is an exact reformulation, not a proof of the inequality.

\section{External Wronskian normal form}

Let
\[
 f(z)=\int_{\mathbb R}K(t)e^{-izt}\,dt.
\]
Writing $t=a+b$ and $s=b$ gives
\[
 \widehat\nu_2(z)
 =\iint_{\mathbb R^2}(a-b)^2K(a)K(b)e^{-iz(a+b)}\,da\,db.
\]
Using
\[
 \int tK(t)e^{-izt}\,dt=i f'(z),
 \qquad
 \int t^2K(t)e^{-izt}\,dt=-f''(z),
\]
one obtains
\[
 \boxed{\widehat\nu_2(z)=2\bigl(f'(z)^2-f(z)f''(z)\bigr).}
\]
Multiplication by $\cosh(ty)$ averages the transforms at $x+iy$ and $x-iy$. Since $\Psi_y(2u)=4D_y(u)$,
\[
 \boxed{
 4\widehat D_y(2x)
 =\Re\!\left(f'(x+iy)^2-f(x+iy)f''(x+iy)\right).
 }
\]
Thus the direct external Fourier target is equivalently
\[
 \boxed{
 \Re\!\left(f'(x+iy)^2-f(x+iy)f''(x+iy)\right)>0
 }
\]
for every real $x$ and every $0<|y|<1/2$.

\section{Complete-monotonicity sufficient route}

Define
\[
 \boxed{H_y(q):=D_y(\sqrt q)=\cosh(2y\sqrt q)C(\sqrt q),\qquad q\ge0.}
\]
Suppose
\[
 \boxed{(-1)^mH_y^{(m)}(q)\ge0\quad\forall m\ge0,\ q\ge0,\quad 0<|y|<\frac12.}
 \tag{CM}
\]
Then Bernstein's theorem provides a positive measure $\mu_y$ with
\[
 H_y(q)=\int_0^\infty e^{-\lambda q}\,d\mu_y(\lambda).
\]
Thus
\[
 D_y(u)=\int_0^\infty e^{-\lambda u^2}\,d\mu_y(\lambda).
\]
The Riemann correlation kernel is integrable and tends to zero at infinity, so a non-zero constant component is excluded. Fourier transformation therefore gives the positive Gaussian mixture
\[
 \widehat D_y(x)=\int_{(0,\infty)}\sqrt{\frac{\pi}{\lambda}}\exp\!\left(-\frac{x^2}{4\lambda}\right)d\mu_y(\lambda)>0.
\]
Consequently
\[
 \boxed{(CM)\text{ for all }0<|y|<\frac12\Longrightarrow RH.}
\]
The implication is exact. Full complete monotonicity remains open, but its first derivative inequality is proved below.

\section{The distinct internal Jensen tilt}

Define instead
\[
 \boxed{J_y(u)=\int_{\mathbb R}r^2\cosh(2yr)K(u+r)K(u-r)\,dr.}
\]
A direct two-variable change of variables gives
\[
 \boxed{4\widehat J_y(2x)=|f'(x+iy)|^2-\Re\!\left(f''(x+iy)\overline{f(x+iy)}\right).}
\]
This is the internal complex Jensen/Laguerre functional.

The two tilted families and their Fourier forms are not identified:
\[
 \boxed{D_y(u)=\cosh(2yu)C(u)\neq J_y(u)}
\]
for general $y\ne0$. They coincide only at
\[
 \boxed{D_0=J_0=C.}
\]
This distinction is part of the proof firewall.

\section{Quantitative curvature already contained in SOH-G004}

SOH-G004 proves uniformly for every theta channel
\[
 g_n''<-12
\]
and separately proves
\[
 \operatorname{Var}_p(g_n')<2.
\]
Its exact log-sum identity therefore yields
\[
 (\log\Phi)''
 =\sum_np_ng_n''+\operatorname{Var}_p(g_n')
 <-12+2=-10.
\]
Hence the full even kernel satisfies
\[
 \boxed{-(\log K)''>10.}
\]
Put $L=-\log K$. Then $L$ is even and
\[
 \boxed{L''>10.}
\]
This quantitative margin is an internal consequence of the existing canonical G004 estimates; no external numerical curvature estimate is used here.

\section{SOH-G024 first-order complete-monotonicity theorem}

Differentiate the centered correlation:
\[
 C'(u)
 =-\int_{\mathbb R}r^2K(u+r)K(u-r)
 \bigl[L'(u+r)+L'(u-r)\bigr]dr.
\]
For $u>0$, evenness gives $L'(u-r)=-L'(r-u)$. Since $L''>10$, the derivative $L'$ is strongly increasing and therefore
\[
 L'(u+r)-L'(r-u)>10\bigl[(u+r)-(r-u)\bigr]=20u.
\]
Consequently
\[
 \boxed{C'(u)<-20uC(u),\qquad u>0,}
\]
or equivalently
\[
 \boxed{-\frac{C'(u)}{C(u)}>20u.}
\]

Now set $u=\sqrt q$. From $D_y(u)=\cosh(2yu)C(u)$,
\[
 -\frac{H_y'(q)}{H_y(q)}
 =-\frac1{2u}\frac{D_y'(u)}{D_y(u)}
 >10-\frac{|y|}{u}\tanh(2|y|u).
\]
Since $\tanh v<v$ for $v>0$,
\[
 \boxed{
 -\frac{H_y'(q)}{H_y(q)}
 >10-2y^2
 >\frac{19}{2}
 }
\]
for every $q>0$ and $0<|y|<1/2$. The continuous limit at $q=0$ obeys the corresponding endpoint inequality. Therefore
\[
 \boxed{H_y'(q)<0}
\]
globally throughout the full Dimitrov--Xu strip.

\begin{theorem}[SOH-G024 first-order complete-monotonicity theorem]
For the actual Riemann kernel and every $0<|y|<1/2$, the radial-square external profile $H_y$ satisfies the first non-trivial complete-monotonicity inequality globally. More precisely,
\[
 -\frac{H_y'(q)}{H_y(q)}>\frac{19}{2}
 \qquad(q>0).
\]
Higher complete-monotonicity orders are not asserted.
\end{theorem}

\section{Gaussian domination}

Integrating the logarithmic-slope estimate gives
\[
 \boxed{H_y(q)<H_y(0)e^{-19q/2}\qquad(q>0).}
\]
Equivalently,
\[
 \boxed{D_y(u)<D_y(0)e^{-19u^2/2}\qquad(u\ne0).}
\]
Thus every external Dimitrov--Xu tilt in the open strip has a uniform Gaussian envelope.

\section{Second-order frontier}

Define
\[
 S_y(q):=-\frac{d}{dq}\log H_y(q).
\]
The preceding theorem proves $S_y(q)>19/2$. Direct differentiation yields
\[
 \boxed{
 \frac{H_y''(q)}{H_y(q)}=S_y(q)^2-S_y'(q).
 }
\]
Therefore
\[
 \boxed{H_y''(q)\ge0\iff S_y'(q)\le S_y(q)^2.}
\]
This Riccati-type inequality is the next exact analytic target. The first-order lower bound alone does not prove it.

\section{Finite diagnostic}

The accompanying receipt evaluates the repository Riemann kernel on a declared finite grid in $q$ and $y$ and applies symmetric five-point finite differences through order four to $H_y(q)$. The first-order sign is independently proved analytically; sampled signs at orders two through four remain classified only as \texttt{FINITE\_DIAGNOSTIC\_NOT\_PROOF}.

\gap{The inequalities $(-1)^mH_y^{(m)}(q)\ge0$ for orders $m\ge2$, all $q\ge0$, and all $0<|y|<1/2$ are not proved. Therefore strict Fourier positivity, Wiener density for the full Riemann family, SOH-G003, SOH-C005, PF$_3$, PF$_\infty$, and RH remain open.}
